\documentclass{article}
\usepackage{graphicx}
\usepackage{authblk}
\usepackage{booktabs}
\usepackage{float}
\usepackage[brazil]{babel}
\usepackage[utf8]{inputenc}
\usepackage[T1]{fontenc}
\usepackage{indentfirst}
\begin{document}
\begin{titlepage}
	\begin{center}
	
	\vspace{115pt}
    \textbf{\Huge{Projeto Programação Concorrente e Paralela - CUDA}}\\
        
	\vspace{115pt}
    Carlos Eduardo Nogueira Silva \\
    Felipe Gomes da Silva \\
    Gabriel Martins Brum \\
    Luis Henrique Salomão Lobato \\
	\end{center}
	
	\vspace{1cm}
	\begin{center}
		\vspace{\fill}
    \large{Novembro, 2025} 
	\end{center}
\end{titlepage}
\date{November 2025}




\section{Introdução}

O mapeamento da propagação de doenças é um tema de grande relevância, principalmente após a pandemia de COVID-19. Através do conhecimento de como uma doença se transmite em uma população, é possível prever seu comportamento e tomar medidas preventivas. 

Neste trabalho, é possível simular como que uma doença seria propagada com base na suas taxas de infecção, recuperação e mortalidade. Os casos de estudo são em mapas bidimensionais, com as seguintes dimensões: 20x20, 200x200 e 2000x2000.

Neste trabalho foram medidos os tempos de execução tanto em uma execução serial (CPU.c) quanto em uma execução paralela (GPU.c) utilizando CUDA.

\section{Metodologia}
\subsection{Cálculo}

\subsection{Testes}

\section{Resultados}

\section{Conclusão}

\end{document}